\section{Literature Review}
\label{sec:literature_review}

TCSSOP jointly decides (i)~how many physical samples to allocate to each project, (ii)~which environmental configuration each climate chamber should take, and (iii)~how to schedule the resulting multi-stage tests. The closest studies therefore fall into four groups that match these decisions: R\&D test scheduling with prototype or sample counts; thermal-chamber and oven scheduling with a shared processing environment; reconfigurable-machine scheduling; and unrelated parallel-machine scheduling with eligibility. We review each group by stating what it contributes to the refrigerator setting and what it still leaves open. Table~\ref{tab:literature_gap} summarizes the comparison.

The studies that come closest on sample and prototype decisions arise in automotive R\&D testing. \citet{BartelsZimmermann2009} schedule a large set of automotive tests on expensive experimental vehicles in order to minimize the number of vehicles required. They formulate the problem as a multi-mode resource-constrained project scheduling model with renewable and cumulative resources, minimum and maximum time lags, and test dependencies, and they solve small instances by mixed-integer linear programming and large instances by priority-rule heuristics. The industrial message is directly relevant: in R\&D validation, the number of physical units is not only a cost item but an operational decision that determines whether a test portfolio can be completed under time pressure. \citet{ShiReichEpelmanKlampflCohn2017} address prototype-vehicle test planning at Ford. Tests must be assigned to compatible vehicles under destructive-test restrictions, time windows, delivery dates, and sequencing constraints. Fit-and-Swap heuristics and integer-programming components replace a manual planning process whose motivation---not leaving expensive units idle, and not building unnecessary ones---is close to the sample-budget versus tardiness trade-off in this study. \citet{LimtanyakulSchwiegelshohn2012} extend this line methodologically. They minimize the number of prototypes through a hybrid master--slave decomposition: a mixed-integer planning model selects how many prototypes and which variants are needed, while a constraint-programming model assigns and sequences tests on those units. That split is conceptually close to the proposed framework, which generates a sample vector, evaluates it by a feasible scheduler, and then updates chamber configurations.

These three papers treat the physical unit as a scarce bin into which a given test set is packed. The test list is exogenous: which tests must be performed is known in advance, and the optimization is to cover that list with as few vehicles as possible, or to pack tests tightly onto a limited fleet. Compatibility is defined by vehicle--test features. In TCSSOP the sample count changes the job set itself. Pull-down is repeated for every physical sample, so an extra sample creates upstream work; at the same time, downstream tests can be distributed across samples and run in parallel. More samples are therefore not always better capacity. Congestion can increase upstream while slack appears downstream, and the net effect on tardiness is non-monotone. The objective also differs. Automotive studies minimize prototype count as the primary goal; this study minimizes total tardiness first and uses total sample usage only lexicographically. Eligibility is determined by chamber temperature--humidity settings, voltage capability, and station capacity, not by vehicle features.

Reliability demonstration testing supplies the complementary sample-size literature. \citet{Fernandez2022} jointly optimizes sample size, acceptance number, and test duration under a Weibull lifetime model and limited expected producer and consumer risks. Sample size is an explicit decision, and the industrial intuition of a trade-off between the number of expensive units and test cost is related to the pull-down cost versus parallelism benefit in TCSSOP. The model is statistical rather than operational. It does not schedule multiple projects on shared chambers, does not encode product-specific test portfolios or stage-based precedence, and does not measure total tardiness. Chamber eligibility, station capacity, sample exclusivity, and environmental configuration are absent. The value of a sample in TCSSOP is therefore not a reliability-risk threshold but the tardiness of the induced feasible schedule.

The closest analog of the chamber as a resource is thermal-oven and thermal-chamber scheduling, where jobs share one processing environment. \citet{LacknerMrkvickaMusliuWalkiewiczWinter2023} introduce the Oven Scheduling Problem in electronics manufacturing. Jobs may occupy the same batch only if their requirements are compatible. Oven eligibility and availability, capacity, release dates, setups between batches, and tardiness are modeled together, and because ovens are energy-intensive the objective also includes cumulative batch processing time. The shared-environment logic is the same as in a climate chamber: all stations inherit one temperature--humidity setting, so incompatible tests cannot share a chamber. Long thermal occupations also match the 5--15 day tests in refrigerator R\&D. \citet{MajumderLahaSuganthan2018} are closer still in resource type. Their application is burn-in and thermal-cycling chambers for assembled printed circuit boards. Jobs with unequal ready times and non-identical sizes are batched on parallel thermal chambers; a hybrid cuckoo search with variable neighborhood search minimizes makespan. The paper is a direct reference for thermal test chambers, long chamber occupancy, and neighborhood-based search.

Neither oven paper decides sample size. Compatibility is a batch attribute or processing-time window, whereas in TCSSOP it is produced by a tactical temperature--humidity configuration and by product-level voltage requirements, and the configuration is held fixed over the horizon. Batch packing groups existing jobs into a common environment; it does not create new pull-down jobs or open extra downstream parallel channels. \citet{LacknerMrkvickaMusliuWalkiewiczWinter2023} may emphasize energy or batch duration; \citet{MajumderLahaSuganthan2018} minimize makespan. This study uses total tardiness as the primary objective and total sample consumption as the secondary one, under a heterogeneous multi-stage R\&D precedence structure that oven batching does not encode.

Chamber configuration is closer to reconfigurable manufacturing. \citet{MahmoodjanlooTavakkoliMoghaddamBaboliBozorgiAmiri2020} add reconfigurable machine tools to the flexible job shop: each operation is eligible for at least one machine--configuration pair, configuration-dependent setups arise when a machine switches state, and the objective is makespan. Configuration and sequencing cannot be separated, which is also true of chamber environments. \citet{FanZhangLiuShenGao2022} restrict auxiliary modules that give a machine its functions. Limited modules create resource conflict when several machines request the same capability, analogous to humidity control and voltage compatibility being available only on a subset of chambers. \citet{FanZhangShenGao2023} combine lot-streaming with machine reconfiguration in a flexible job shop and minimize total weighted tardiness with a genetic-algorithm matheuristic and mixed-integer local search on lot sizes. Quantity-related decisions then affect tardiness in the same model as reconfiguration, which is the job-shop interaction closest to TCSSOP. \citet{VahediNouriRohaninejadHanzalekFoumani2025} assign orders to eligible machine--configuration pairs on non-identical parallel reconfigurable machines, then batch and schedule under capacity limits. The decision is not only which machine, but which machine in which configuration---the same logic as assigning a climate chamber to a temperature--humidity setting.

These models still differ on three points. First, configurations typically switch between consecutive operations, with setups paid at each change. In the industrial laboratory, frequent reconfiguration is undesirable because of stabilization and interruption of ongoing tests, so environments are fixed over the planning horizon and updated only when sample-driven workload changes. Second, lot-streaming is not sample-size optimization. Lot-streaming splits a given production quantity into sublots; it does not create jobs. An extra refrigerator sample can increase total processing effort through repeated pull-down while also enabling downstream parallelism, and that trade-off is non-monotone. Third, the setting is manufacturing or remanufacturing orders, not a multi-project R\&D test portfolio on shared climate chambers with voltage restrictions and sample exclusivity.

After chamber environments are fixed, assignment and sequencing resemble unrelated parallel-machine scheduling with eligibility. \citet{MaeckerShenMonch2023} minimize total weighted tardiness on unrelated machines with eligibility and delivery times, which is the standard language for heterogeneous chambers that cannot process every test. \citet{PerezGonzalezFernandezViagasZamoraGarciaFraminan2019} develop constructive heuristics for unrelated machines with eligibility and setup times, supporting a light earliest-due-date decoder inside repeated search. \citet{AfzaliradRezaeian2016} combine resource constraints, sequence-dependent setups, precedence, and machine eligibility in one model---the unrelated-machine constraint package closest to station capacity, environmental eligibility, and stage-based precedence. \citet{SantoroJunqueira2023} add machine availability to eligibility in a recent formulation. In all four studies the job set is exogenous. Jobs, processing times, and eligible-machine sets are data; choosing a sample count does not change the number of operations or their duration profile. Machines are also independent. In TCSSOP every station in a chamber shares one environment, so eligibility is decided at chamber-configuration level rather than only at machine--job level. Sample exclusivity---a physical unit cannot run two tests at once---is a more specific R\&D restriction than a generic additional resource. Unrelated-machine methods therefore guide the scheduling layer, but they do not resolve sample optimization or chamber reconfiguration.

\begin{table}[H]
\centering
\caption{Closest related studies and the remaining gap relative to TCSSOP}
\label{tab:literature_gap}
\footnotesize
\setlength{\tabcolsep}{3.2pt}
\renewcommand{\arraystretch}{1.12}
\begin{tabular}{@{}>{\raggedright\arraybackslash}p{3.35cm}>{\raggedright\arraybackslash}p{2.25cm}>{\raggedright\arraybackslash}p{2.45cm}>{\raggedright\arraybackslash}p{2.55cm}>{\raggedright\arraybackslash}p{2.25cm}@{}}
\toprule
\textbf{Study} & \textbf{Application} & \textbf{Configurable resource} & \textbf{Sample/prototype decision} & \textbf{Shared environment} \\
\midrule
\citet{BartelsZimmermann2009}; \citet{ShiReichEpelmanKlampflCohn2017}; \citet{LimtanyakulSchwiegelshohn2012} & Automotive R\&D tests & Prototype vehicle as the machine & Cover a given test set with a fleet & No \\
\citet{Fernandez2022} & Reliability demonstration tests & Not modeled & Statistical sample size and duration & No \\
\citet{LacknerMrkvickaMusliuWalkiewiczWinter2023}; \citet{MajumderLahaSuganthan2018} & Thermal ovens and cycling chambers & Recipe or batch environment & No; jobs are given & Yes, as a batch \\
\citet{MahmoodjanlooTavakkoliMoghaddamBaboliBozorgiAmiri2020}; \citet{FanZhangLiuShenGao2022}; \citet{FanZhangShenGao2023}; \citet{VahediNouriRohaninejadHanzalekFoumani2025} & Reconfigurable manufacturing & Machine, module, or line configuration & Lot-streaming at most; no sample-generated jobs & No \\
\citet{AfzaliradRezaeian2016}; \citet{MaeckerShenMonch2023}; \citet{PerezGonzalezFernandezViagasZamoraGarciaFraminan2019}; \citet{SantoroJunqueira2023} & Unrelated parallel machines & Eligibility, not multi-feature reconfiguration & No & No \\
This study (TCSSOP) & Refrigerator climate chambers & Temperature, humidity, and voltage & Endogenous sample size & Yes, at chamber level \\
\bottomrule
\end{tabular}
\end{table}

Table~\ref{tab:literature_gap} makes the remaining gap precise. No existing study jointly (i)~treats sample size as an endogenous generator of pull-down workload and of downstream parallelism, (ii)~configures chambers that have multiple reconfigurable features and share one environment across stations, and (iii)~schedules the resulting multi-stage refrigerator tests under a tardiness objective. That combination is the problem addressed by TCSSOP.

To close this gap, this study formulates TCSSOP and develops an iterative framework that coordinates chamber-environment assignment, sample-size optimization, and feasible schedule construction. Unlike conventional neighborhood search that perturbs sequencing or machine assignment, the proposed diversification mechanism perturbs sample sizes and thereby explores alternative workload structures.
